\section{AI Decision Layer}\label{sec:decision}

\subsection{Role}

The optional decision layer adds a quality judgement to the detector's class label. It uses
THX-02~\cite{HALX2026thx02}, a decision model that answers typed questions over text and images
with probabilities, through its hosted API. The layer never commands motion: it only fills
decision fields of a part record, and the PLC acts on them inside its own interlocks and
bin-selection rule (Section~\ref{sec:plc}). The PLC does not depend on it: with
the layer switched off, \texttt{GVL\_Bridge.xDecisionActive} is FALSE and the PLC decides each
part by its rule as soon as vision has identified it.

\subsection{Questions and Bin Selection}

One request per part carries the part's image crop, the detected class and confidence, and two
questions: a \emph{choice} between \texttt{ok}, \texttt{rework} and \texttt{scrap} (each defined by a
short criterion), and a yes/no question whether the identification is reliable. The client maps
the answers to \texttt{ST\_DecisionMsg}: the chosen quality with its probability, and
\emph{unreliable} when the probability of ``no'' is at least 0.8. In \texttt{F\_SelectBin} an
unreliable identification sends the part to Quarantine, Scrap and Rework go to their bins, and
only parts judged \texttt{ok} reach their type bin. The detector's class and the decision both
change where the robot places the part.

\subsection{Timing, Rate Limit and Fallback}

The decision for a part must arrive within the PLC's deadline of 2.5\,s after the capture; the
client's time budget is that deadline minus the time already spent. Requests run in a worker
thread with its own event loop, and the bridge polls the answers once per 20\,ms cycle, so a slow
or stalled API cannot delay the cycle. The free tier of the API counts every question, so a part
costs two decisions; a client-side token bucket allows 50 decisions per minute (burst~2), i.e.\
25 parts per minute, below the 60 per minute per address that the API allows. A request that
would exceed the deadline, a rate-limit answer (HTTP~429) that asks for more than 1\,s of waiting,
or any other error ends in the fallback. Fallback decisions are never written to the PLC: the
part keeps waiting until the PLC's own deadline passes and is then decided by the rule, with a
\emph{DecisionFallback} warning on the HMI. A decision that arrives after its part has reached
the pick window is dropped.

\subsection{Record and Replay}

Every exchange is written to a JSON-lines file (run id, seed, part id, time stamp, latency,
request hash, response or fallback reason). A replay mode feeds a recording back into a run, so
results can be recomputed without the API, whose hosted model is not version-pinned.

\subsection{Status}

The decision layer, its rate limiter, deadline handling and record/replay are implemented and
covered by unit tests with a mocked HTTP API, and the request format was checked with one call
using a synthetic image. The layer was \emph{not} used in any of the cell runs reported in
Section~\ref{sec:evaluation}; in those runs the PLC decided every part by its rule. The twin
also does not yet generate defective parts, so the quality decision has no ground truth, and the
planned comparison of the rule baseline with the decision layer is future work. The alarm-triage
advice in the original design was not implemented.
